State Hilbert's basis theorem and give the main idea of its proof

If \(R\) is commutative and noethern, so is \(R[x]\).

Take any ideal \(\mathfrak{a} \in R[x]\), consider the ideal \(I\) consisting of all leading coeff.
of the polynomials in \(\mathfrak{a}\), as \(R\) is noethern, I is fin. gen. choosing a fin. gen.
we can construct an ideal \(\mathfrak{a}' \subset \mathfrak{a}\) of polynomials, which have exactly 
these \(a \in I\) as leading coefficients (which is again fin. gen.). One can then show that \(\mathfrak{a} = (\mathfrak{a} \cap M) + \mathfrak{a}'\),
where \(M \subset R[x]\) is the submodule generated by the monomials \(\{1, x, \dots, x^{r-1}\), as both \((\mathfrak{a} \cap M)\) and \(\mathfrak{a}\)
are fin. gen. so must \(\mathfrak{a}\), hence it is noetherian.


Moreover it follows that \(R[x_1, \dots, x_n]\) is noetherian under the same assumptions.

Let

Define simple/irreducible modules and their relations to idempotent elements

A \(R\)-module \(M\) is simple iff. \(\{0\}, M\) are the only submodules of \(M\).
Equivalently if \(M = M_1 \oplus M_2\) then \(M_i = 0\) for one of them.

An example are vector spaces of dimension 1.

Every integral domain is simple (where the submodules are ideals).

Let \(e \in \text{End}_{R}(M)\) be idempotent, then \(e(M)\) and \((1-e)(M)\) are submodules
of \(M\), satisfying \(M = e(M) \oplus (1 - e)(M)\).
The proof is a verification.

Moreover an \(R\)-module \(M \neq \{0\}\) is simple if and only if \(0, 1\) are the only idempotent
elements in \(\text{End}_R(M)\).

Proof: Use the previous for the implication and for the converse note that for \(\pi : M \to M_1\), \(\iota : M_2 \to M\)
\(e = \pi \circ \iota\) defines an idempotent element.

Can noetherian modules be decomposed into simple modules?

Yes indeed, the idea is to take the set of all submodules \(\Sigma\) which can not be decomposed
into a direct sum of simple modules and show that its empty. For contradiction assume it isn't so,
in particular then \(M \in \Sigma\). Let \(N \in \chi\) be the set of all proper submodules, for which a 
\(M' \in \Sigma\) exists, s.t. \(M = N \oplus M'\), this is non-empty since \(\emptyset \in \chi\), by assumption.

As \(M\) is noetherian there is a maximum \(N \in \chi\) and corresponding \(M' \in \Sigma\), then \(M'\) must be decomposable
in at least one non-simple module \(M' = M_1 \oplus M_2\), so one of them is also in \(\Sigma\), w.l.o.g. let it be \(M_1\) 
hence \(M_2 \oplus N\) is another member of \(\chi\) since\(M = M_1 \oplus (M_2 \oplus N)\), contradicting \(N\)'s maximality.

State and proof the rank and basis properties for submodules of free \(M\) over a PID. 

Let \(R\) be a principal domain, let \(M\) be a free \(R\)-module of finite rank \(r\) and \(N \subset M\) a submodule. 
Then
1. \(N\) is free of rank \(r' \leq r\)
2. There is a basis \(\{x_n\}_{n \in [r]}\) of \(M\) s.t. \(\{a_nx_n\}_{n \in [r']}\) is a basis of \(N\) where \(a_1, \dots, a_{r'} \in R\)
satisfying \(a_1 | a_2 | \dots | a_{r'}\).


Note the degenerate case \(N = \{0\}\) is already free. Henceforth consider for non-trivial \(N\), \(\Sigma = \{(a_{\phi}) | \phi \in \text{Hom}_R(M, R)\}\), where \(a_{\phi} = \phi(N)\),
(\(\phi(N)\) is a submodule of \(R\) so its an ideal and it must have a generator as \(R\) is PID). Since \(R\) is also noetherian, there is some maximum
\(a_1\) in \(\Sigma\) with corresponding \(\nu(y) = a_1, y \in N\). We show that \(a_1 | \phi(y)\) for all \(\phi \in \text{Hom}_R(M, R)\), let \(d = r_1 a_1 + r_2 \phi(y)\)
be the principal generator of the ideal \((\phi(y) \cap a_1\) (thus by Bezout's theorem we have the previous equation). Let \(\psi = r_1\nu + r_2\phi : M \to R\) then 
\(\phi(y) = d\), thus \(d \in \psi(N)\) and \( (d) \subseteq \psi(N)\). But \(d\) is a divisor of \(a_1\) so we have also \((a_1) \subset (d)\) but by maximality of (a_1)
we must in fact have equality \((a_1) = (d) = (\psi(N))\) and by assumption \(a_1 | d | \phi(y)\).

Consider now \(\pi_i(x) = b_i : M \to R\), for \(x \in M\) and \(x = \sum_{i = 1}^{r} b_i x_i\), hence \(a_1 | \pi_i(y)\) for all \(i\). 
Thus we can write \(\pi_i(y) = a_1 c_i\) and define \(y_1 = \sum_{i = 1}^{r} b_i x_i\). Note that \(a_1y_1 = y\). Since \(a_1 = \nu(y) = \nu(a_1y_1) = a_1 \nu(y_1)\) we must have \(\nu(y_1) = 1\)
(\(R\) is an integral domain). We then assert

3. \(M = Ry_1 \oplus \text{ker}(\nu)\)
4. \(N = Ra_1y_1 \oplus (N \cap \text{ker}(\nu))\)

To show 3., take any \(x \in M\) and write \(x = \nu(x)y_1 + (x - \nu(x)y_1)\), since 

\(\nu(x - \nu(x)y_1) = \nu(x) - \nu(x)\nu(y_1) = \nu(x) - \nu(x) 1 = 0\), \(x - \nu(x)y_1 \in \text{ker}(\nu)\). To show that the sum is direct, suppose \(ry_1 \in \text{ker}(\nu)\) then clearly \(r = 0\).
For 4. observe that \(a_1 | \nu(x')\) for all \(x' \in N\), thus \(\nu(x') = ca_1\) for some \(c \in R\) and \(x' = \nu(x')y_1 + (x' - \nu(x'))y_1 = ca_1y_1 + (x' - ca_1y_1)\), we have \(x' - ca_1y_1 \in \text{ker}(\nu)\) and directness follows the same way as for 3.
We proceed via induction (the base case \(\text{rk}(N) = 0\) is already satisfied).
Suppose \(\text{rk}(N) = r'\) then \(\text{rk}(N \cap \text{ker}(\nu)) = r' - 1\), by induction hypothesis this module is free and adjoining \(a_1y_1\) gives a basis of \(N\), so \(N\) is also free, which was 1.
For 2. take 3. and apply the induction hypothesis again to see that \(\text{ker}(\nu)\) is free and must have a basis \(\{x_n\}_{n \in 2 \dots r}\). Then from 4. there is \(\{a_nx_n\}_{n \in 2 \dots r'}\)
a basis of \(\text{ker}(\nu) \cap N\) and even more, we immediately know that \(N\) is free with decomposition \(\{a_nx_n\}_{n \in 1 \dots r'}\). It remains to show that \(a_1 | a_2\), since this is not part of 4 (the remaining divisors are by induction hypothesis). Construct the hom. \(\phi(y_1) = \phi(y_2) = 1, \phi(y_i) = 0 \text{otherwise}\), then \(\phi(a_1y_1) = a_1\) and \((a_1) = \phi(N)\) by maximality of \(a_1\), but also \(\phi(a_2y_2) = a_2\), hence \(a_2 \in \phi(N) = (a_1)\) giving the desired result.


with suitable divisors


Suppose \(\text{rk}(N) = r'\) then \(\text{rk}(N \cap \text{ker}(\nu)) = r' - 1\), by induction hypothesis this module is free and adjoining \(a_1y_1\) gives a basis of \(N\), so \(N\) is also free.



Proof and state the existence of the invariant factor form

Let \(R\) be a principal domain, let \(M\) be a fin. gen. \(R\)-module.
1. \(M\) is isomorphic to the direct sum of finitely many cyclic modules. More precisely
\(M \simeq R^{r} \oplus R/(a_1) \oplus R/(a_2) \dots \oplus R/(a_m)\)
where \(a_i \in R\) are not units and satisfy \(a_1 | a_2 | \dots | a_m\).
2. \(M\) is torsion free iff. \(M\) is free
3. \(Tor(M) = R/(a_1) \oplus R/(a_2) \dots \oplus R/(a_m)\)


Proof:
Let \(\{x_k\}_{k \in [n]}\) be a set of generators for \(M\) of minimal cardinality. Let \(R^{n}\) be the free \(R\)-module 
with basis \(\{b_k\}_{k \in [n]}\) and the surjective \(\pi : R^{n} \to M, \pi(b_k) = x_k\). By the first isomorphism theorem
\(R^{n}/\text{ker}(\pi) \simeq M\). By the theorem of bases for submodules of free modules there is a basis \(\{a_ky_k\}_{k \in [m]}\)
of \(\text{ker}(\pi)\) satisfying \(a_1 | a_2 | \dots\). Note that \(\phi: Ry_1 \oplus Ry_2 \dots Ry_n \to R/(a_1) \oplus R/(a_2) \dots \oplus R^{n-m}\)
\(\phi(\alpha_1 y_1, \dots \alpha_n y_n) = (\alpha_1 \text{mod} (a_1), \alpha_2 \text{mod} (a_2), \dots, \alpha_{m+1}, \dots, \alpha_n \) has \(\text{ker}(\phi) = Ra_1y_1 \oplus Ra_2y_2 \dots \oplus Ra_my_m\). 
Hence \(M \simeq Ry_1 \oplus \dots Ry_m/Ra_1y_1 \oplus Ra_2y_2 \dots \oplus Ra_my_m \simeq \text{im}(\phi) \simeq R/(a_1) \oplus R/(a_2) \dots \oplus R^{n-m} \).

The rest follows naturally.

It can be shown that this decomposition is actually unique up to multiplication by units for the invariant factors \(a_i\).

Proof and state the existence of the elementary divisors form

Let \(R\) be a principal domain, let \(M\) be a fin. gen. \(R\)-module.
Then \(M \simeq R^{r} \oplus R/(p^{\alpha_1}_1) \oplus R/(p^{\alpha_2}_2) \oplus \dots \oplus R/(p^{\alpha_t}_t)\)
where \(\alpha_i > 0\) and \(p_i\) are prime in \(R\).

Proof:
We know there exists a invariant factor form for \(M\), moreover any PID \(R\) is also UFD, so for each \(a_i = u p_1^{\alpha_1}\dots p_s^{\alpha_s}\)
It holds \((p_i^{\alpha_i}) + (p_j^{\alpha_j}) = R\), for \(i \neq j\), so they're comaximal pairs. Their intersection in turn is \((a_i)\) (as the \(\text{lcm}(p_1^{\alpha_1}\dots p_s^{\alpha_s})\))
By the chinese remainder theorem \(R/(a_i) = R/(p_1^{\alpha_1}) \oplus R/(p_2^{\alpha_2}) \oplus \dots \oplus R/(p_s^{\alpha_s})\). Doing this for each factor yields the elementary divisor form.


It can be shown that this decomposition is actually unique up to multiplication by units for the elementary divisors \(p_i^{\alpha_i}\).


State the relation between \(M\) and \(pM\) for a prime \(p \in R\), \(R\) a PID.

Denote the field \(F = R/(p)\), since \((p)\) is maximal in \(R\).
It holds
1. If \(M = R^n\), then \(M/pM \simeq F^n\)
2. If \(M = R/(a)\), where \(a \neq 0\) in \(R\), \(M/pM \simeq F\) if \(p | a\) otherwise \(M/pM \simeq \{0\}\)
3. Let \(M = \bigoplus_{a_i} R/(a_i)\) with \(p | a_i\) for all \(i\), then \(M/pM \simeq F^n\).
